% HD-G17. GPT-6 Astra Pro, inherited effort, /root/g34.
% Candidate data: Severn, Hexagon2610.00477v1 / Zenodo23549514.
% No theorem or saved rank output from that preprint is a premise.
\subsection{HD-G17: a smooth rank-nine point from an independent Laurent calculation}
\label{sec:HD-G17}

The actual twelve-point cover constructed below has primitive
Gaussian multiplicities $(3,3)$ and period differential of rank $9$.
Its Hermitian determinant class is $[-3]$, by the direct cup
calculation HD-S23 on the exact tuple displayed here.  Thus this component
dominates the corresponding nine-dimensional Weil family.
The numerical inputs are candidate residue centres from
B.~Severn's unrefereed
\href{https://doi.org/10.5281/zenodo.23549514}{certificate archive};
the existence verification is the fresh HD-I17 proof and the
multiplication calculation here was independently rebuilt.

\paragraph{1. The exact system and the actual oriented tuple.}
Let $a,b,c$ be the three integral matrices in the independently checked
$E_8$ group certificate, in the fundamental-weight basis, and put
$d=(abc)^{-1}$.  HD-I17 constructs an actual Fuchsian system
\[
 dY_M=\left(\frac{B_0}{z}+\frac{B_1}{z-1}
                  +\frac{B_c}{z-c_0}\right)Y_M\,dz,\qquad c_0=2+i,
\]
where
\[
 B_0=\operatorname{diag}(1,3,7,9,11,13,17,19)/20,\qquad
 B_1=u_1v_1^t,\quad B_c=u_2v_2^t,
 \quad v_1^tu_1=v_2^tu_2=\tfrac12.
\]
The exact monodromies at the three finite poles, for the based
polygonal loops used there, are simultaneously conjugate to
$a,b,c$.  Its positive infinity monodromy is $d$.
This uses the actual orientation calculation in HD-I17: from
$z_0=1/2-i$ the stem directions have counterclockwise order
$c_0,1,0$, so a positive large finite boundary traverses the three
loops in that time order; transport for $Y'=AY$ multiplies in
the reverse order, giving $abc$.

On the $s$-line take
\begin{equation}
             Y(s)=s^{\,1/4-3}Y_M(s^5).                 \tag{G17.1}
\end{equation}
The integer part of the scalar gauge has trivial monodromy.
The finite branch points are $0$, the roots of $s^5=1$ and $s^5=c_0$,
with infinity the twelfth point.  Schreier representatives for
$z=s^5$ give the following positively oriented tuple:
\begin{equation}
 \begin{aligned}
 \bigl(&ia^5,\ a^{-4}ba^4,a^{-4}ca^4,\
        a^{-3}ba^3,a^{-3}ca^3,\\
       &a^{-2}ba^2,a^{-2}ca^2,\
        a^{-1}ba,a^{-1}ca,\ b,c,\ -id^5\bigr).
 \end{aligned}                                      \tag{G17.2}
\end{equation}
Indeed, the lifted loops around a root of $1$ or $c_0$ are the
displayed conjugates, the loop at zero lifts five turns of the
zero loop, and the loop at infinity lifts five turns of the infinity
loop.  The scalar factor has monodromy $i$ at zero, $-i$ at infinity,
and $1$ at the other ten points.  The word identity
\[
 a^5\prod_{j=4,3,2,1,0}a^{-j}(bc)a^j=(abc)^5
\]
proves that (G17.2) has product $1$, including its scalar factors.
These words and all twelve matrices are recorded in
\path{work/g34/e8/actual_pullback_tuple.json}.

Conjugating its first eleven generators by $a^4$ gives the set
\[
          \{ia^5\}\ \cup\
          \{a^kba^{-k},a^kca^{-k}:0\le k\le4\}.
\]
The independent finite group certificate HD-S22 therefore gives the full
group $G=W(E_8)\cdot\mu_4$.  Its local codimensions are
$4,1,\ldots,1,8$; the ages are $2,\tfrac12,\ldots,\tfrac12,4$.
Consequently both primitive Hodge multiplicities equal $11-8=3$.
The direct cup calculation HD-S23 on (G17.2) gives determinant $-24$,
whose norm class is $[-3]$ because $8=N_{\mathbf Q(i)/\mathbf Q}(2+2i)$.
No inverse--transpose comparison is used in this argument.
For comparison of supplier class data only, the last entry is
conjugate to the candidate archive's last entry by $a^{-1}$:
$a^{-1}(-id^5)a=-ia^{-1}d^5a$; all ten other nonscalar entries
are reflection conjugates.

\paragraph{2. The local analytic gauges needed for this calculation.}
The residue $B_\infty=-B_0-B_1-B_c$ has the exact eigenvalue multiset
\[
 \beta_\infty=
 \left(-1,-1,-\frac{11}{12},-\frac23,-\frac7{12},
                       -\frac5{12},-\frac13,-\frac1{12}\right).
                                                        \tag{G17.3}
\]
HD-I17 proves these labels using outward-rounded Gershgorin bounds
and the exact finite infinity monodromy.  Let $P$ be an eigenbasis
in the displayed order.

Here the original-system local gauges require no resonant
normal-form assertion.  At zero all eigenvalue differences of $B_0$
have absolute value below $1$, and the same holds for (G17.3).
For a system $dF/dz=(A_0/z+A_{\rm reg}(z))F$, substitution
$F=H(z)z^{A_0}$, $H(0)=1$, gives
\[
 (m-\operatorname{ad}A_0)H_m
       =\sum_{\ell=0}^{m-1}(A_{\rm reg})_\ell H_{m-1-\ell}.
\]
Each operator is invertible for $m\ge1$.  In an eigenbasis its
inverse entries are bounded by a constant divided by $m$;
the usual scalar majorant for this recurrence converges on a
small disc, proving that $H$ is holomorphic.
One may first use the same recurrence with Jordan blocks:
its denominators still avoid zero.  The finite local monodromy
then makes $A_0$ semisimple, since a nonzero Jordan nilpotent would
give a nontrivial unipotent factor in $\exp(2\pi iA_0)$.
This also proves the needed diagonalizability of $B_\infty$.
At the two reflection poles the eigenvalues are $0$ and $1/2$,
so the same argument applies there.

After (G17.1) the gauges at zero and infinity are series in $s^5$
and $t^5$, respectively, where $t=1/s$.  Their constant terms are
the identity in the zero frame and $P$ in the infinity frame.
The exact new exponents are
\begin{equation}
 \begin{aligned}
 \lambda_0&=
 \left(-\frac52,-2,-1,-\frac12,0,\frac12,\frac32,2\right),
 \\
 \lambda_\infty&=
 \left(-\frac94,-\frac94,-\frac{11}6,-\frac7{12},
                         -\frac16,\frac23,\frac{13}{12},\frac73\right).
 \end{aligned}                                          \tag{G17.4}
\end{equation}
Their floors are
\begin{equation}
 \begin{aligned}
 \lfloor\lambda_0\rfloor&=(-3,-2,-1,-1,0,0,1,2),\\
 \lfloor\lambda_\infty\rfloor&=(-3,-3,-2,-1,-1,0,1,2).
 \end{aligned}                                          \tag{G17.5}
\end{equation}

\paragraph{3. Actual differential spaces and their complete Laurent description.}
Let $C$ be the smooth compactification of the connected finite cover
defined by (G17.2).  On it $Y$ is single-valued away from the
ramification points and obeys $g^*Y=Y\rho(g)$.
Thus
\[
       \omega=Y^{-1}r(s)\,ds,\qquad
       \nu=Y^tq(s)\,ds                              \tag{G17.6}
\]
identify rational vectors $r,q$ satisfying holomorphy with the
multiplicity spaces for the two inverse representations.
Explicitly, the
components of $\omega$ transform by $\rho(g)^{-1}$, and hence
specify a map from the dual representation; the components of
$\nu$ specify a map from the representation itself.
Conversely such equivariant holomorphic forms give $G$-invariant
$Y\omega/ds$ and $Y^{-t}\nu/ds$, which descend to rational vectors
on the sphere.  Meromorphic extension at the finitely many poles
follows from the local power-times-holomorphic gauges just proved.
This proves both directions of the identification.

At a branch point with exponent $\lambda$ and inertia order $e$,
write $s-p=y^e$.  A coefficient of base order $k$ in the first
expression in (G17.6) has upstairs order
$ek-e\lambda+e-1$.  Since $e\lambda$ is integral, this is nonnegative
exactly when $k\ge\lfloor\lambda\rfloor$.
For the second expression the condition is
$k\ge\lfloor-\lambda\rfloor$.

For $r$, the reflection-point exponents $0,1/2$ have floor zero,
so $r$ is regular at all ten of those points.  The minimum floors
at zero and infinity give the complete expression
\begin{equation}
                   r(s)=\sum_{k=-3}^{1} r_k s^k,
                   \qquad r_k\in\mathbf C^8.          \tag{G17.7}
\end{equation}
At infinity the vector multiplying $dt$ is
$-t^{-2}r(1/t)$.  Hence its coefficient at order $-2-k$ is $-r_k$.
The complete linear constraints on (G17.7) are
\begin{align}
 (r_k)_j&=0 &&\text{for }k<\lfloor\lambda_{0,j}\rfloor,
                                                        \tag{G17.8}\\
 (P^{-1}r_k)_j&=0
       &&\text{for }-2-k<\lfloor\lambda_{\infty,j}\rfloor .
                                                        \tag{G17.9}
\end{align}
There are $20+17=37$ rows on $40$ unknowns.
No omitted gauge coefficient affects them: at zero the lowest
power is $-3$, so the first possible gauge correction has order $2$;
at infinity the same statement holds.  Every imposed condition
has order at most $1$, by (G17.5).

For $q$, at a root of $s^5=1$ the coefficient of its possible
simple pole must lie in $\mathbf Cv_1$.  Indeed the seven
zero-eigenvectors of $B_1$ span $\ker(v_1^t)$, and their annihilator
is precisely $\mathbf Cv_1$.  The corresponding direction at a
root of $s^5=c_0$ is $\mathbf Cv_2$.
Interpolation at the five distinct roots therefore gives the
complete rational expression
\begin{equation}
 q(s)=\sum_{k=-2}^{1}q_ks^k+
       v_1\frac{p_1(s)}{s^5-1}
          +v_2\frac{p_c(s)}{s^5-c_0},
 \qquad \deg p_1,\deg p_c\le4.                       \tag{G17.10}
\end{equation}
The remainder after removing these principal parts has only its
zero and infinity poles, which proves completeness and uniqueness.
There are $4\cdot8+5+5=42$ coefficients.

The zero constraints are the coordinate vanishing conditions of
orders $k<\lfloor-\lambda_{0,j}\rfloor$.
For $k=0,1$ the coefficient to use is
\[
        q_k-v_1[p_1]_k-c_0^{-1}v_2[p_c]_k;
\]
for $k=-2,-1$ it is $q_k$.  These give $16$ rows.
At infinity, after the common minus sign from $ds=-t^{-2}dt$,
the coefficient at an order $\ell\le1$ is
\[
 q_{-\ell-2}
        +v_1[p_1]_{3-\ell}+v_2[p_c]_{3-\ell},          \tag{G17.11}
\]
with coefficients outside their specified index ranges set to zero.
Apply $P^t$ and require its $j$th coordinate to vanish for
$\ell<\lfloor-\lambda_{\infty,j}\rfloor$.
These give $23$ more rows, hence a $39$ by $42$ matrix.
Again the first corrections from the original local gauges occur
in order at least $3$ at zero and $2$ at infinity; every required
condition has order at most $1$.  Thus (G17.8)--(G17.11) are the
full holomorphy conditions, not truncated approximations.

\paragraph{4. The exact product and the independent interval computation.}
The invariant contraction of the actual forms (G17.6) is
\begin{equation}
                    \nu^t\omega=q(s)^tr(s)\,ds^2.     \tag{G17.12}
\end{equation}
It has at most simple poles downstairs at the twelve branch points.
For example, an invariant upstairs holomorphic quadratic differential
in $y$ has the form $y^a(dy)^2$ with $a\ge0$ and
$a+2\equiv0\pmod e$; expressing it in $s=y^e$ gives downstairs
order $(a+2)/e-2\ge-1$.
Consequently multiplication by
$s(s^5-1)(s^5-c_0)$ identifies (G17.12) with a polynomial of degree
at most $8$.  Conversely all such polynomials give the full
nine-dimensional space
$H^0(\mathbf P^1,K_{\mathbf P^1}^2(B))$.

The fresh program
\path{work/g34/e8/independent_laurent_rank.py}
implements precisely the two matrices just described and the nine
polynomial products.  It does not import the supplied period
matrices, Frobenius coefficients, pivots or rank outputs.
Its only numerical inputs are the $87$ residue/conjugator centre
coordinates, the fixed dyadic gauge data, and the enclosing radius
\[
                   \rho_L=1.0528417230446196\cdot10^{-34}.
\]
The primitive residue unknowns are contained in this box by the
fresh HD-I17 computation: its exact-root displacement is bounded
above by $1.2859780343191978\cdot10^{-35}<\rho_L$.
The JSON conversion was made only after inspection of the pickle
opcodes and use of a restricted loader allowing the sole global
$\mathtt{builtins.complex}$; the resulting data schemas and hashes
are recorded.

All arithmetic after pivot selection uses Arb through
python-flint~0.9.0 at $256$ bits.  Floating-point Gaussian
elimination with complete pivoting selects row and column sets only.
For each selected square block $S$, a ball inverse encloses its
inverse and a point matrix $R$ satisfies
$\|1-RS\|_\infty<1$; this certifies invertibility for every exact
matrix enclosed.  The free-coordinate solution then encloses an
actual kernel basis.  The seven infinity eigenspaces are obtained
the same way at their exact rational eigenvalues, of multiplicities
two and one; their combined matrix $P$ is certified invertible.
The fresh infinity Gershgorin check also lies within $10^{-10}$
of these exact labels.

For the $37$ by $40$ and $39$ by $42$ matrices the respective
Neumann bounds are less than
$5.965\cdot10^{-29}$ and $2.572\cdot10^{-30}$.
Thus both kernels have dimension exactly $3$.
The coefficient matrix of their nine products in the basis
$1,s,\ldots,s^8$ has determinant enclosed by
\begin{align*}
 \Re(\det)&=-2.40174170406026598352216
                         \ \pm 4.29\cdot10^{-24},\\
 \Im(\det)&=\phantom{-}0.188389348347788894425049
                         \ \pm 8.76\cdot10^{-25}.
\end{align*}
The certified absolute lower bound exceeds $2.409$.
The matrix entries, pivot indices, all bounds, input hashes and
fresh log are in
\path{independent_laurent_rank.json}
and \path{independent_laurent_rank.log} beside the program.
The unused polynomial coefficients outside degrees $0$ through $8$
enclose zero; this is a consistency check, whereas their exact
vanishing was proved geometrically in (G17.12).

\paragraph{5. Identification with the actual period differential.}
Infinitesimal deformation of the cover is deformation of the
twelve marked quotient points, since the finite monodromy datum
has no infinitesimal deformation.  Its tangent space is
$H^1(\mathbf P^1,T_{\mathbf P^1}(-B))$, dual by Serre duality
to the nine-dimensional quadratic space above.
On the curve, the derivative of the Hodge filtration is cup
product with the Kodaira--Spencer class.  Serre duality makes its
transpose multiplication of holomorphic one-forms.

Restrict this identity to the two inverse representation factors.
The invariant map from their two multiplicity spaces into
the symmetric square of $H^0(K_C)$ is the actual coevaluation
contraction
\[
       f\otimes h\longmapsto
               \sum_{j=1}^{8} f(e_j^*)h(e_j).
\]
Under (G17.6) this is exactly (G17.12).  More explicitly, a tangent
map of the multiplicity Hodge structure acts on its isotypical
factor as $1_{\mathbf C^8}\otimes u$.  The trace pairing evaluates
the displayed coevaluation on this map as $8$ times its evaluation
on $u$, because $\operatorname{Tr}(1_{\mathbf C^8})=8$.
Thus the cotangent identification differs by the explicit invertible
scalar $8$; its rank agrees with (G17.12).
Taking the finite quotient trace multiplies the descended quadratic
differential by the nonzero degree $|G|$; therefore it has the same
rank.  The actual primitive abelian factor HD-S22 and its polarized
isogeny comparison with this multiplicity Hodge structure were
constructed in the separate group/factor argument.  Their maps on
weight-one cohomology are invertible, so this rank is precisely the
period rank of that factor.

The computed rank is $9$, equal to both the Hurwitz dimension
$12-3$ and the dimension of the signature-$(3,3)$ unitary domain.
The holomorphic submersion theorem therefore gives a nonempty
open period image.  This proves dominance for the actual component
and the determinant class calculated directly from (G17.2).

\paragraph{Verification scope.}
The author read the complete proof and all new Laurent code on
21 November 2026.  The supplied L program was also fully read
and freshly rerun after safe data conversion, giving its separate
nonzero residue minor; that rerun is not represented as an
independently written algorithm.  The proof here uses the distinct
Laurent reconstruction above.  The independent M verification,
the actual-tuple cup calculation, and this new rank computation
are separately recorded events, with their exact data links.
Nothing in this statement supplies algebraic cycles: the supplier
construction and its group-to-row verification are separate parts
of the overall argument.
